\hypertarget{ejercicio-4.-autuxf3matas-finitos-en-jflap}{%
\section{Ejercicio 4. Autómatas finitos en JFLAP}\label{ejercicio-4.-autuxf3matas-finitos-en-jflap}}

\hypertarget{muxe9todo-de-diseuxf1o}{%
\subsection{Método de diseño}\label{muxe9todo-de-diseuxf1o}}

Todos los autómatas se diseñaron con la misma pregunta: \emph{¿qué es lo mínimo que hay que recordar del prefijo leído para decidir el resto?} La respuesta define los estados; la función de transición dice cómo cambia ese recuerdo al leer un símbolo, y los estados de aceptación son aquellos cuyo recuerdo ya garantiza que la cadena pertenece al lenguaje. Cuatro familias de enunciados cubren casi todas las listas:

\begin{longtable}[]{@{}
  >{\raggedright\arraybackslash}p{(\columnwidth - 4\tabcolsep) * \real{0.1667}}
  >{\raggedright\arraybackslash}p{(\columnwidth - 4\tabcolsep) * \real{0.6667}}
  >{\raggedright\arraybackslash}p{(\columnwidth - 4\tabcolsep) * \real{0.1667}}@{}}
\caption{Familias de diseño.}\tabularnewline
\toprule\noalign{}
\begin{minipage}[b]{\linewidth}\raggedright
Familia
\end{minipage} & \begin{minipage}[b]{\linewidth}\raggedright
Qué guarda cada estado
\end{minipage} & \begin{minipage}[b]{\linewidth}\raggedright
Ejemplos
\end{minipage} \\
\midrule\noalign{}
\endfirsthead
\toprule\noalign{}
\begin{minipage}[b]{\linewidth}\raggedright
Familia
\end{minipage} & \begin{minipage}[b]{\linewidth}\raggedright
Qué guarda cada estado
\end{minipage} & \begin{minipage}[b]{\linewidth}\raggedright
Ejemplos
\end{minipage} \\
\midrule\noalign{}
\endhead
\bottomrule\noalign{}
\endlastfoot
Conteo módulo k & El residuo de una longitud, de un número de apariciones o del valor numérico leído & 097, 139, 145, 147 \\
Búsqueda de patrones & El sufijo más largo leído que es prefijo del patrón (construcción de Knuth-Morris-Pratt) & 099, 102, 103, 106 \\
Producto de autómatas & Un par de recuerdos, uno por condición, cuando el enunciado pide «y» o «o» & 107, 109, 133 \\
Prefijos y extremos & Cuánto del prefijo exigido se ha leído, o el primer símbolo y el último & 120, 135, 137 \\
\end{longtable}

Cuando una condición ya no puede cumplirse (por ejemplo, la cadena no empieza como se exige) se usa un estado \textbf{sumidero}: no es de aceptación y todas sus transiciones vuelven a él. En los enunciados de la Lista 3 que piden un AFN, el AFD se obtuvo convirtiendo ese AFN y minimizando el resultado; su ficha indica qué estados del AFN representa cada estado del AFD.

\hypertarget{implementaciuxf3n-y-validaciuxf3n}{%
\subsection{Implementación y validación}\label{implementaciuxf3n-y-validaciuxf3n}}

Cada autómata se capturó en JFLAP 7.1 y se exportó a \texttt{automatas/lista2/} y \texttt{automatas/lista3/}; los archivos \texttt{.jff} se leen también con el software del Ejercicio 5 de la Práctica 2. Se resolvieron los 77 ejercicios con asterisco que piden un autómata; los ejercicios 116 (lenguaje no regular), 143 (figuras ausentes en la lista), 191 y 201--235 no piden un AFD y se resuelven en las tareas.

Cada ficha del Anexo A tiene el diagrama en JFLAP, la tabla de transiciones ---en todos los autómatas, no sólo en los cinco que pide el enunciado---, tres cadenas válidas, tres inválidas y los casos frontera: la cadena vacía, la aceptada más corta y, cuando existe, una cadena justo por debajo de esa longitud. Las cadenas se eligieron de longitudes crecientes y se comprobaron dos veces: con \emph{Input → Multiple Run} en JFLAP y con un programa que compara cada autómata contra la definición del lenguaje sobre todas las cadenas de longitud ≤ 14 (≤ 9 con tres símbolos). El archivo \texttt{automatas/pruebas.json} reúne las cadenas de todas las fichas.

\hypertarget{simulaciuxf3n-paso-a-paso}{%
\subsection{Simulación paso a paso}\label{simulaciuxf3n-paso-a-paso}}

La simulación paso a paso (\emph{Input → Step with Closure}) se usó en los cinco autómatas de la tabla; cada captura muestra la configuración (estado y resto de la entrada) después de algunos pasos y el resultado final.

\begin{longtable}[]{@{}
  >{\raggedright\arraybackslash}p{(\columnwidth - 6\tabcolsep) * \real{0.1579}}
  >{\raggedright\arraybackslash}p{(\columnwidth - 6\tabcolsep) * \real{0.5263}}
  >{\raggedright\arraybackslash}p{(\columnwidth - 6\tabcolsep) * \real{0.1579}}
  >{\raggedright\arraybackslash}p{(\columnwidth - 6\tabcolsep) * \real{0.1579}}@{}}
\caption{Autómatas simulados paso a paso.}\tabularnewline
\toprule\noalign{}
\begin{minipage}[b]{\linewidth}\raggedright
Ejercicio
\end{minipage} & \begin{minipage}[b]{\linewidth}\raggedright
Lenguaje
\end{minipage} & \begin{minipage}[b]{\linewidth}\raggedright
Cadena simulada
\end{minipage} & \begin{minipage}[b]{\linewidth}\raggedright
Resultado
\end{minipage} \\
\midrule\noalign{}
\endfirsthead
\toprule\noalign{}
\begin{minipage}[b]{\linewidth}\raggedright
Ejercicio
\end{minipage} & \begin{minipage}[b]{\linewidth}\raggedright
Lenguaje
\end{minipage} & \begin{minipage}[b]{\linewidth}\raggedright
Cadena simulada
\end{minipage} & \begin{minipage}[b]{\linewidth}\raggedright
Resultado
\end{minipage} \\
\midrule\noalign{}
\endhead
\bottomrule\noalign{}
\endlastfoot
097 & Longitud múltiplo de 3 & \texttt{102211} & aceptada en r0 \\
113 & Dos apariciones de 010, con traslape & \texttt{01010} & aceptada \\
120 & Primer y último símbolo distintos & \texttt{2101} & aceptada \\
135 & bᵐabⁿ con m, n \textgreater{} 0 & \texttt{bbab} & aceptada \\
147 & Binario (izquierda a derecha) múltiplo de 6 & \texttt{1100} & aceptada (12 = 2·6) \\
\end{longtable}

\evidencia[0.95]{../evidencias/jflap/lista2-097-paso-a-paso.png}{JFLAP, simulación paso a paso de \texttt{102211} en lista2-097.}

\evidencia[0.95]{../evidencias/jflap/lista2-113-paso-a-paso.png}{JFLAP, simulación paso a paso de \texttt{01010} en lista2-113.}

\evidencia[0.95]{../evidencias/jflap/lista2-120-paso-a-paso.png}{JFLAP, simulación paso a paso de \texttt{2101} en lista2-120.}

\evidencia[0.95]{../evidencias/jflap/lista2-135-paso-a-paso.png}{JFLAP, simulación paso a paso de \texttt{bbab} en lista2-135.}

\evidencia[0.95]{../evidencias/jflap/lista2-147-paso-a-paso.png}{JFLAP, simulación paso a paso de \texttt{1100} en lista2-147.}

\hypertarget{resumen-de-los-autuxf3matas}{%
\subsection{Resumen de los autómatas}\label{resumen-de-los-autuxf3matas}}

\begin{longtable}[]{@{}lllcl@{}}
\caption{Autómatas de los ejercicios con asterisco.}\tabularnewline
\toprule\noalign{}
Ejercicio & Σ & Diseño original & Estados del AFD & Archivo \\
\midrule\noalign{}
\endfirsthead
\toprule\noalign{}
Ejercicio & Σ & Diseño original & Estados del AFD & Archivo \\
\midrule\noalign{}
\endhead
\bottomrule\noalign{}
\endlastfoot
L2 · 097 & \{0, 1, 2\} & AFD & 3 & \texttt{lista2-097.jff} \\
L2 · 099 & \{0, 1, 2\} & AFD & 7 & \texttt{lista2-099.jff} \\
L2 · 102 & \{a, b\} & AFD & 7 & \texttt{lista2-102.jff} \\
L2 · 103 & \{a, b\} & AFD & 7 & \texttt{lista2-103.jff} \\
L2 · 106 & \{0, 1, 2\} & AFD & 8 & \texttt{lista2-106.jff} \\
L2 · 107 & \{0, 1, 2\} & AFD & 13 & \texttt{lista2-107.jff} \\
L2 · 109 & \{0, 1, a, b, c\} & AFD & 22 & \texttt{lista2-109.jff} \\
L2 · 111 & \{0, 1, 2\} & AFD & 21 & \texttt{lista2-111.jff} \\
L2 · 112 & \{0, 1, 2\} & AFD & 7 & \texttt{lista2-112.jff} \\
L2 · 113 & \{0, 1, 2\} & AFD & 7 & \texttt{lista2-113.jff} \\
L2 · 118 & \{a, b\} & AFD & 6 & \texttt{lista2-118.jff} \\
L2 · 120 & \{0, 1, 2\} & AFD & 7 & \texttt{lista2-120.jff} \\
L2 · 121 & \{0, 1, 2\} & AFD & 6 & \texttt{lista2-121.jff} \\
L2 · 122 & \{0, 1, 2\} & AFD & 4 & \texttt{lista2-122.jff} \\
L2 · 123 & \{a, b\} & AFD & 15 & \texttt{lista2-123.jff} \\
L2 · 124 & \{0, 1\} & AFD & 13 & \texttt{lista2-124.jff} \\
L2 · 125 & \{0, 1\} & AFD & 2 & \texttt{lista2-125.jff} \\
L2 · 126 & \{0, 1\} & AFD & 2 & \texttt{lista2-126.jff} \\
L2 · 127 & \{0, 1\} & AFD & 2 & \texttt{lista2-127.jff} \\
L2 · 128 & \{0, 1\} & AFD & 6 & \texttt{lista2-128.jff} \\
L2 · 131 & \{0, 1\} & AFD & 2 & \texttt{lista2-131.jff} \\
L2 · 133 & \{a, b\} & AFD & 4 & \texttt{lista2-133.jff} \\
L2 · 135 & \{a, b\} & AFD & 5 & \texttt{lista2-135.jff} \\
L2 · 137 & \{a, b\} & AFD & 4 & \texttt{lista2-137.jff} \\
L2 · 139 & \{0, 1\} & AFD & 2 & \texttt{lista2-139.jff} \\
L2 · 141 & \{0, 1\} & AFD & 3 & \texttt{lista2-141.jff} \\
L2 · 145 & \{0, 1\} & AFD & 6 & \texttt{lista2-145.jff} \\
L2 · 147 & \{0, 1\} & AFD & 7 & \texttt{lista2-147.jff} \\
L2 · 148 & \{A, N, P, S\} & AFD & 13 & \texttt{lista2-148.jff} \\
L2 · 149 & \{0, 1\} & AFD & 3 & \texttt{lista2-149.jff} \\
L2 · 150 & \{0, 1, a\} & AFD & 4 & \texttt{lista2-150.jff} \\
L3 · 151 & \{0, 1\} & AFD & 10 & \texttt{lista3-151.jff} \\
L3 · 152 & \{0, 1, 2\} & AFD & 6 & \texttt{lista3-152.jff} \\
L3 · 153 & \{0, 1\} & AFD & 7 & \texttt{lista3-153.jff} \\
L3 · 154 & \{0, 1\} & AFD & 11 & \texttt{lista3-154.jff} \\
L3 · 155 & \{0, 1\} & AFD & 12 & \texttt{lista3-155.jff} \\
L3 · 156 & \{0, 1\} & AFD & 5 & \texttt{lista3-156.jff} \\
L3 · 157 & \{0, 1\} & AFD & 6 & \texttt{lista3-157.jff} \\
L3 · 158 & \{0, 1\} & AFD & 3 & \texttt{lista3-158.jff} \\
L3 · 159 & \{0, 1\} & AFD & 3 & \texttt{lista3-159.jff} \\
L3 · 160 & \{0, 1\} & AFD & 3 & \texttt{lista3-160.jff} \\
L3 · 161 & \{0, 1\} & AFD & 3 & \texttt{lista3-161.jff} \\
L3 · AFN del 097 & \{0, 1, 2\} & AFD & 3 & \texttt{lista3-afn-097.jff} \\
L3 · AFN del 099 & \{0, 1, 2\} & AFN & 7 & \texttt{lista3-afn-099.jff} \\
L3 · AFN del 102 & \{a, b\} & AFN & 7 & \texttt{lista3-afn-102.jff} \\
L3 · AFN del 103 & \{a, b\} & AFN & 7 & \texttt{lista3-afn-103.jff} \\
L3 · AFN del 106 & \{0, 1, 2\} & AFN & 8 & \texttt{lista3-afn-106.jff} \\
L3 · AFN del 107 & \{0, 1, 2\} & AFN & 13 & \texttt{lista3-afn-107.jff} \\
L3 · AFN del 109 & \{0, 1, a, b, c\} & AFN & 22 & \texttt{lista3-afn-109.jff} \\
L3 · AFN del 111 & \{0, 1, 2\} & AFD & 20 & \texttt{lista3-afn-111.jff} \\
L3 · AFN del 112 & \{0, 1, 2\} & AFN & 7 & \texttt{lista3-afn-112.jff} \\
L3 · AFN del 113 & \{0, 1, 2\} & AFN & 7 & \texttt{lista3-afn-113.jff} \\
L3 · AFN del 118 & \{a, b\} & AFN & 5 & \texttt{lista3-afn-118.jff} \\
L3 · AFN del 120 & \{0, 1, 2\} & AFN & 7 & \texttt{lista3-afn-120.jff} \\
L3 · AFN del 121 & \{0, 1, 2\} & AFD & 5 & \texttt{lista3-afn-121.jff} \\
L3 · AFN del 122 & \{0, 1, 2\} & AFD & 3 & \texttt{lista3-afn-122.jff} \\
L3 · AFN del 123 & \{a, b\} & AFN & 15 & \texttt{lista3-afn-123.jff} \\
L3 · 189 & \{0, 1\} & AFN & 7 & \texttt{lista3-189.jff} \\
L3 · 190 & \{0, 1\} & AFN & 19 & \texttt{lista3-190.jff} \\
L3 · 192 & \{0, 1\} & AFN & 5 & \texttt{lista3-192.jff} \\
L3 · 193 & \{0, 1\} & AFN & 6 & \texttt{lista3-193.jff} \\
L3 · 194 & \{0, 1, 2\} & AFN & 2 & \texttt{lista3-194.jff} \\
L3 · 195 & \{0, 1\} & AFN & 7 & \texttt{lista3-195.jff} \\
L3 · 236 & \{0, 1\} & AFN-λ & 4 & \texttt{lista3-236.jff} \\
L3 · 237 & \{0, 1\} & AFN-λ & 4 & \texttt{lista3-237.jff} \\
L3 · 238 & \{0, 1\} & AFN-λ & 6 & \texttt{lista3-238.jff} \\
L3 · 239 & \{0, 1\} & AFN-λ & 6 & \texttt{lista3-239.jff} \\
L3 · 240 & \{0, 1\} & AFN-λ & 6 & \texttt{lista3-240.jff} \\
L3 · 254 & \{0, 1\} & AFD & 2 & \texttt{lista3-254.jff} \\
L3 · 255 & \{0, 1\} & AFD & 2 & \texttt{lista3-255.jff} \\
L3 · 256 & \{0, 1, 2\} & AFD & 2 & \texttt{lista3-256.jff} \\
L3 · 257 & \{0, 1\} & AFD & 3 & \texttt{lista3-257.jff} \\
L3 · 258 & \{0, 1\} & AFD & 3 & \texttt{lista3-258.jff} \\
L3 · 260 & \{0, 1\} & AFN & 6 & \texttt{lista3-260.jff} \\
L3 · 261 & \{0, 1\} & AFN & 3 & \texttt{lista3-261.jff} \\
L3 · 262 & \{a, b\} & AFD & 5 & \texttt{lista3-262.jff} \\
L3 · 263 & \{0, 1\} & AFN & 3 & \texttt{lista3-263.jff} \\
\end{longtable}
